\section{附录}\label{sec:appendix}

\subsection{符号说明}

\begin{table}[htbp]
    \centering
    \caption{本文使用的符号}
    \label{tab:notation}
    \begin{tabular}{@{}ll@{}}
        \toprule
        \textbf{符号} & \textbf{含义} \\
        \midrule
        $\mathbb{N}$ & 自然数集合 \\
        $\mathbb{P}$ & 素数集合 \\
        $\pi(x)$ & 不超过 $x$ 的素数个数 \\
        $\Lambda(n)$ & 冯·曼戈尔特函数 \\
        $\mu(n)$ & 莫比乌斯函数 \\
        $\phi(n)$ & 欧拉函数 \\
        $e(x)$ & $e^{2\pi i x}$ \\
        $f \ll g$ & $f = O(g)$ \\
        $f \sim g$ & $\lim_{x \to \infty} f(x)/g(x) = 1$ \\
        $P_r$ & 素因子个数不超过 $r$ 的整数 \\
        $\mathfrak{S}(n)$ & 奇异级数 \\
        \bottomrule
    \end{tabular}
\end{table}

\subsection{重要定理证明概要}

\subsubsection{维诺格拉多夫三素数定理证明概要}

\begin{proof}[证明概要]
    设 $r_3(n)$ 表示奇数 $n$ 表示为三个素数之和的方法数。使用圆法：
    \begin{equation}
        r_3(n) = \int_0^1 S(\alpha)^3 e(-n\alpha) \, d\alpha
    \end{equation}
    
    将积分分为优弧 $\mathfrak{M}$ 和劣弧 $\mathfrak{m}$：
    \begin{equation}
        r_3(n) = \int_{\mathfrak{M}} S(\alpha)^3 e(-n\alpha) \, d\alpha + \int_{\mathfrak{m}} S(\alpha)^3 e(-n\alpha) \, d\alpha
    \end{equation}
    
    优弧贡献主要项：
    \begin{equation}
        \int_{\mathfrak{M}} S(\alpha)^3 e(-n\alpha) \, d\alpha \sim \frac{n^2}{2(\log n)^3} \mathfrak{S}(n)
    \end{equation}
    
    劣弧贡献低阶项，使用维诺格拉多夫中值定理估计：
    \begin{equation}
        \int_{\mathfrak{m}} S(\alpha)^3 e(-n\alpha) \, d\alpha = o\left(\frac{n^2}{(\log n)^3}\right)
    \end{equation}
\end{proof}

\subsubsection{陈氏定理证明概要}

\begin{proof}[证明概要]
    陈景润的证明综合使用了筛法和解析数论技术。主要步骤包括：
    
    \begin{enumerate}
        \item 构造加权筛法函数
        \item 应用布伦-蒂奇马什定理
        \item 使用维诺格拉多夫定理估计指数和
        \item 通过精细分析得到最终结果
    \end{enumerate}
    
    具体地，陈景润证明了：
    \begin{equation}
        \sum_{\substack{p \leq x \\ x-p = P_2}} 1 \gg \frac{x}{(\log x)^2}
    \end{equation}
    其中 $P_2$ 表示不超过两个素数乘积的数。
\end{proof}

\subsection{历史文献}

\begin{enumerate}
    \item Goldbach, C. (1742). Letter to L. Euler, June 7, 1742.
    \item Hardy, G. H., \& Littlewood, J. E. (1923). Some problems of 'Partitio numerorum'; III: On the expression of a number as a sum of primes. \textit{Acta Mathematica}, 44, 1-70.
    \item Vinogradov, I. M. (1937). Representation of an odd number as a sum of three primes. \textit{Doklady Akademii Nauk SSSR}, 15, 291-294.
    \item Chen, J. R. (1966). On the representation of a larger even integer as the sum of a prime and the product of at most two primes. \textit{Kexue Tongbao}, 17, 385-386.
    \item Helfgott, H. A. (2013). The ternary Goldbach conjecture is true. \textit{arXiv preprint arXiv:1312.2719}.
\end{enumerate}
